\section{Introduction}
\label{sec:introduction}

Recent clinical developments \cite{bib1} with the therapeutic agents based on Ac-225 prompted a vigorous search for practical methods to produce this isotope at scale. The clinical value of this isotope is estimated so high, that current market price of 1 patient dose of the isotoep alone currently goes as high as \$10,000 USD, higher than any other isotope ever tried in humans. This high value started a frenzy of development and over a dozen of approaches were suggested employing many types of accelerators, nuclear reactors and isotopes generators. However, practical considerations of production at scale, with the right purity, in a distributed manner and robust enough to satisfy pharmaceutical manufacturing constrains eliminate most of the suggested methodologies. 3 methods currently supply most of the market: Th-229/Ac-225 generator, Ra-226(p,2n)Ac-225 reaction enabled by low energy cyclotron and Ra-226(\gamma,2n)Ac-225 process enabled by high energy eLINACs.
The last methodology is rapidly gaining popularity. The reason for this popularity is not the great advantages of the process, but rather lack of significant problems typically associated with the competitors. A common (p,2n) approach produces product contaminated with Ac-227; additionally potential failure of the radium target in this process can irreversibly contaminate the entire accelerator. Th-229/Ac-225 generator potentially is a very attractive technology, but the starting material is available commercially to only one company in the world, making this approach of no interest for any other developer. Photonuclear route also comes with significant problems. First, radium starting material is scantly available and hard to deal with in routine production. Second, eLINACs of sufficient power are unique instruments with less than a dozen of suitable units built in the entire world ever. Third, low conversion efficiency requires chemical processing of curies of Ra-226 to obtain meaningful quantities of Ac-225, necessitating large and complex GMP hot cell facilities. None of these problems, however are of fundamental nature, they all can be solved with the right investment and a capable team. As the old Jewish saying goes, "if a problem can be solved with money, it isn't a problem, it's an expense".
This review aims to provide a chemists perspective to the photonuclear production route. 

